\chapter{FRONTEND TECHNOLOGY STACK}

\section{Overview of Technology Stack}

The frontend of the Bus Ticket Booking System is developed using modern web technologies that support interactive user interfaces, responsiveness, and efficient client-side processing. The chosen technology stack ensures smooth navigation, real-time updates, and an enhanced user experience across various devices. Since the scope of the project primarily focuses on frontend development, the selected tools are optimized for handling user interactions, interface design, and state management effectively.

The project adopts the MERN stack approach, with emphasis on the frontend components. Technologies such as HTML, CSS, JavaScript, and React.js are used to build a dynamic single-page application. These technologies collectively provide a scalable and maintainable structure, enabling easy expansion of features such as seat selection, checkout flow, user authentication, and admin dashboard functionalities.

\section{HTML (HyperText Markup Language)}

HTML is used as the core markup language for structuring the web pages of the Bus Ticket Booking System. It defines the basic layout and organization of the application, including navigation bars, forms, buttons, cards, tables, and textual content.

In this project, HTML elements are utilized to design key interfaces such as user login and registration modals, bus search forms, seat selection layouts, checkout pages, invoice views, and the admin dashboard. Semantic HTML tags are applied wherever possible to improve accessibility and readability. The structured nature of HTML ensures seamless integration with CSS for styling and JavaScript for interactivity.

\section{CSS (Cascading Style Sheets)}

CSS is used to enhance the visual appearance and layout of the application. It controls styling elements such as colors, fonts, spacing, alignment, and overall page structure. The application uses CSS to maintain a consistent design theme across all pages, improving usability and visual coherence.

Responsive design techniques are implemented using CSS to ensure the system functions correctly on desktops, tablets, and mobile devices. Bootstrap is also utilized to speed up layout development and provide pre-designed UI components. Proper use of CSS improves readability, aesthetics, and overall user satisfaction.

\section{JavaScript}

JavaScript plays a crucial role in adding interactivity and client-side logic to the application. It enables dynamic updates, form validations, event handling, and conditional rendering of components without reloading the page.

In the Bus Ticket Booking System, JavaScript is used to manage user inputs, handle search operations, validate passenger details, control login sessions, and simulate backend interactions using stored data. It also helps in handling role-based access, such as restricting checkout functionality to logged-in users and providing special privileges to admin users.

\section{React.js}

React.js is a JavaScript library used for building fast and interactive user interfaces, especially for single-page applications. It allows developers to create reusable UI components and manage application state efficiently using hooks.

In this project, React.js is used to structure the application into modular components such as Home, Search Results, Seat Selection, Checkout, Invoice, Login, Register, and Admin Dashboard. React Router is used for navigation between different pages without full page reloads. The use of React’s virtual DOM improves performance by updating only the required parts of the user interface. Overall, React.js enhances scalability, maintainability, and performance of the Bus Ticket Booking System.
